\subsection{树}

当谈到图的道路或连通分支时，指的是底箭图的道路或连通分支。图的两个顶点属于同一连通分支，当且仅当存在连接它们的道路。

设 G 为图。若 \(c = (a_{0}, f_{1}, a_{1}, \ldots, f_{n}, a_{n})\) 是 G 中的道路，则序列 \(\overline{c} = (a_{n}, \overline{f_{n}}, \ldots, a_{1}, \overline{f_{1}}, a_{0})\) 是 G 中的道路，称为 c 的相反道路。设 c 与 \(c'\) 是 G 中可并列的道路；则 \(\overline{c'}\) 与 \(\overline{c}\) 可并列，且有 \(\overline{c * c'} = \overline{c'} * \overline{c}\) 。

若 c 中没有一对相邻的箭是相反的，则称 G 中的道路 c 是无回溯的。设 \(c = (a_{0}, f_{1}, \ldots, f_{n}, a_{n})\) 为 G 中连接 \(a_{0}\) 与 \(a_{n}\) 的道路。若对满足 \(1 \leqslant i < n\) 的整数 i，箭 \(f_{i}\) 与 \(f_{i+1}\) 相反，则 \((a_{0}, f_{1}, \ldots, a_{i-1}, f_{i+2}, \ldots, f_{n}, a_{n})\) 是 G 中连接 \(a_{0}\) 与 \(a_{n}\) 的道路，其长度严格小于道路 c 的长度。由归纳法，因此存在 G 中连接 \(a_{0}\) 与 \(a_{n}\) 的无回溯道路。

定义 2. —— 森林是这样的图：其中每条无回溯的圈都是常值圈。树是连通的森林。

命题 1. —— 设 G 为图。G 的每个森林都含于 G 的一个极大森林中；特别地，存在 G 的极大森林。

为使 G 的森林是极大的，必须且只需其顶点集等于 G 的顶点集，且其连通分支就是 G 的连通分支。

设 \(A_{0}\) 为 G 的森林。G 的森林所成的集，赋有序关系“A 是 B 的子图”后，是归纳的。因此，存在 G 的极大森林，\(A_{0}\) 是它的子图（E, III, 第 21 页, 推论 1）。

以 \(\operatorname{Som}(G)\) 为顶点集、箭集为空的 G 的子图是 G 的森林。因此存在 G 的极大森林。

设 A 为 G 的极大森林。我们证明 A 与 G 有相同的顶点集。以 \(\operatorname{Som}(G)\) 为顶点集、以 \(\operatorname{Fl}(A)\) 为箭集的 G 的子图是一个森林，且 A 是它的子图。于是有 \(\operatorname{Som}(A) = \operatorname{Som}(G)\) 。

现在证明 A 与 G 有相同的连通分支。由于 A 的每个箭都是 G 的箭，A 的每个连通分支含于 G 的一个连通分支。既然 A 与 G 有相同的顶点集，只需证明 G 中处于同一连通分支的两个顶点也处于 A 的同一连通分支。若不然，关系 \(R_{A}\) （“处于 A 的同一连通分支”）严格细于关系 \(R_{G}\) ，且存在 G 的两个顶点，它们不处于 A 的同一连通分支，但却被 G 的一个箭 f 连接。设 B 为 G 的定向子图，其顶点集为 \(\operatorname{Som}(G)\) ，箭集为 \(\mathrm{Fl(A)} \cup \{f, \overline{f}\}\) 。我们来证明 B 是 G 的森林。设 \(c = (a_{0}, f_{1}, \ldots, f_{n}, a_{n})\) 为 B 中长度最小的非常值无回溯圈。由于 A 是森林，圈 c 不是 A 中的圈。设 i（相应地，j）为 \(\{0, \ldots, n\}\) 中使 \(a_{0}\) 与 \(a_{i}\) （相应地，\(a_{j}\) 与 \(a_{n}\) ）不处于 A 的同一连通分支的最小（相应地，最大）整数。这意味着箭 \(f_{i}\) 与 \(f_{j+1}\) 是 B 中与边 \(\{f, \overline{f}\}\) 伴随的定向箭，且它们相反。由于圈 c 无往返，\(f_{i+1} \neq \overline{f_{i}}\) ，因此 \(i \neq j\) ，而道路 \((a_{i}, f_{i+1}, a_{i+1}, \ldots, f_{j}, a_{j})\) 是 B 中长度 < n 的非常值无回溯圈，与 c 长度最小的假设矛盾。由此推出 B 是森林。这与 A 是 G 的极大森林的假设矛盾。

现设 A 为 G 的森林，满足 \(\mathrm{Som}(\mathrm{A}) = \mathrm{Som}(\mathrm{G})\) 与 \(\pi_{0}(\mathrm{A}) = \pi_{0}(\mathrm{G})\) 。我们来证明它是 G 的极大森林。只需证明：若 \(f \notin \mathrm{Fl}(\mathrm{A})\) ，则 G 的以 \(\mathrm{Som}(\mathrm{G})\) 为顶点集、以 \(\mathrm{Fl}(\mathrm{A}) \cup \{f, \overline{f}\}\) 为箭集的子图 B 不是森林。按假设，点 \(o(f)\) 与 \(t(f)\) 处于 A 的同一连通分支；因此存在 A 中连接 \(o(f)\) 与 \(t(f)\) 的无回溯道路 c。于是道路 \(c * \overline{f}\) 是 B 中无回溯且非常值的圈，这表明 B 不是森林。

推论. —— 连通图的极大森林是它的极大树。

注 1. —— 在 LIE, IV, 第 33 页附录中，组合图的概念定义为偶 (A, S)，其中 S 是一个集合，

而 A 是 \(\mathfrak{P}(S)\) 的由二元集组成的子集；S 的元素称为顶点，A 的元素称为边，且两个顶点 \(x\) 与 \(y \in S\) 称为相连的，如果 \(\{x, y\}\) 是一条边。

对这样的组合图 \(\Gamma = (\mathrm{A},\mathrm{S})\) ，伴随一个图 G，其顶点集为 S，箭集 \(\widetilde{\mathrm{A}}\) 是 \(\mathrm{S}^2\) 中由相连顶点的偶所组成的子集，其源映射与靶映射分别与 \(\mathrm{S}^2\) 到 S 的第一、第二投影相合，而对合 \(f\mapsto \overline{f}\) 由 \(\widetilde{\mathrm{A}}\) 上的映射 \((x,y)\mapsto (y,x)\) （\(\mathrm{S}^2\) 到自身的）的限制给出。把 G 的边 \(\{f,\overline{f}\}\) 映为集合 \(\{o(f),t(f)\}\) 的映射是从 G 的边集到组合图 \(\Gamma\) 的边集上的双射。

反之，每个使得每个箭的源与靶不同、且箭由其源与靶确定的图，都具有这种形式。

读者可以验证，组合图的连通性、树或森林的概念，与与之伴随的图的相应概念一致。

